


Aunque el concepto de internet de la cosas o \gls{iot} por sus siglas en inglés es relativamente nuevo, haciendo su aparición por primera vez entre los años 2008 y 2009\cite{Evans2011TheEverything} y definido según el Internet Business Solutions Group (IBSG) de Cisco como “el punto en el tiempo en el que se conectaron a internet más cosas u objetos que personas”\cite{Alcaraz2014InternetCosas}. Hoy en día el IoT tiene un gran impacto en la vida cotidiana de las personas, a cada momento se están creando dispositivos que tienen conexión a internet, por lo que actualmente no son sólo las personas quienes se conectan a la red, sino que también los objetos necesitan tener una conexión para poder funcionar correctamente o realizar las tareas para los que fueron creados. 

El IoT tiene aplicaciones medioambientales, industriales, en ciudades, en el hogar, de forma personal, etc.  Cuando se habla de aplicaciones medioambientales; se refiere a sensores destinados a la protección y a la salud, no solo del ser humano sino también de nuestro planeta. Como lo es el censado de polución en mares y ríos: midiendo el estado del agua y su PH, salinidad, temperatura, iones, etc. Además existen sistemas centrados en la protección de animales por medio de geolocalización, aplicaciones para las catástrofes, entre otros. Los sensores tradicionalmente instalados que solo podían ser leídos en sus centros de control, ahora son capaces de comunicarse mediante GSM y ser fácilmente accedidos (entre ellos y con las centrales). Un claro ejemplo es el proyecto ALARMS de la BGS (Bristish Geological Survey, 2010) que ha conseguido integrar, en un dispositivo de bajo coste con capacidad GSM, sensores miniaturizados para alertar de manera temprana la inestabilidad del suelo (previendo así los corrimientos de tierra)\cite{Valle2014}.

El medio industrial es en donde más aplicaciones se desarrollan; en la actualidad algunas de las aplicaciones existentes son la monitorización de estructuras, seguridad o análisis de datos (sistemas de big-data) y varios más. También hay aplicaciones en el desarrollo de implementos con tecnología IoT en las ciudades como son en los estacionamientos, contenedores inteligentes y control eléctrico.  

Las aplicaciones en el hogar están referidas en mayor medida al control del medio ambiente dentro de la casa como los sistemas de riego autónomo para plantas. En la categoría de aplicaciones pensadas específicamente para las personas, existen una gran variedad de desarrollos como: píldoras inteligentes, trajes de control para bebés, etc\cite{Valle2014}.

Estas nuevas tecnologías a nuestra disposición pueden ser utilizadas prácticamente en todos los ámbitos de la vida de las personas o de las industrias debido a la gran cantidad de posibilidades que ofrece el IoT para obtener, procesar y transmitir información. Cuando se habla de la maleabilidad de estas tecnologías en los diferentes ámbitos de la vida cotidiana, se resaltan aplicaciones como el control de los electrodomésticos o cualquier otro instrumento del hogar que tenga el potencial de recoger y/o procesar datos, monitorear el estado de la vivienda o incluso verificar nuestra salud\cite{Alcaraz2014InternetCosas}. Sin embargo, las aplicaciones anteriormente mencionadas sólo benefician a los propietarios en sí, pero éstas igualmente pueden ser extendidas a la comunidad al plantear aplicaciones de largo alcance como una red inteligente, luces de la calle inteligentes y autos inteligentes conectados. 

Gracias a la gran cantidad de posibilidades que ofrece el IoT, numerosas compañías han introducido soluciones innovadoras para el mercado, por lo que se han creado diferentes infraestructuras para el manejo y control del IoT, entre algunas de esas infraestructuras se pueden mencionar las siguientes: SigFox, que está concebida como una red de telecomunicaciones de cobertura amplia y centrada en dispositivos de bajo consumo\cite{330ohms2018QueFunciona}; Symphony, que se especializa en superar las dificultades que se presentan en LoPoWAN (Low-Power Wide Area Network) y LoRaWAN, el cual es un esquema para abordar la comunicación de largo alcance también conocido como LoRa\cite{Singh2016ComparativeSecurity}. De igual manera en la que se desarrollaron infraestructuras para el manejo de IoT, se pueden mencionar algunas tecnologías que son ampliamente utilizadas en el despliegue y éxito de los productos y servicios basados en IoT, como pueden ser: identificación por radio frecuencia (RFID), redes de sensores inalámbricas (WSN), middleware, computación en la nube y aplicaciones de software de IoT\cite{Lee2015TheEnterprises}. 

Como se mencionó previamente, LoRa es un sistema de telecomunicaciones inalámbricas de largo alcance, baja potencia y baja tasa de bits, promovido como una solución de infraestructura para el Internet de las cosas. Los dispositivos finales usan LoRa en un salto inalámbrico para comunicarse con el gateway conectado a Internet y éstos actúan como puentes transparentes además de permitir el envío de mensajes de retransmisión entre estos dispositivos finales y un servidor de red central\cite{Augustin2016AThings}. Este sistema de telecomunicaciones promete conexiones bidireccionales seguras, bajo consumo de energía, largo alcance de comunicación, bajas velocidades de datos, baja frecuencia de transmisión, movilidad y servicios de localización en aplicaciones de IoT, al tiempo que mantiene simples las estructuras de red y la administración\cite{LoRaAlliance2015AboutAlliance}. Esta tecnología ha recibido mucha atención en los últimos meses de los operadores de red y proveedores de soluciones. Sin embargo, como toda tecnología, tiene limitaciones que deben entenderse claramente para evitar expectativas infladas y desilusión.  

Por ejemplo, si es requerido para aplicaciones en donde sólo se transmite un paquete corto por día, un gateway puede llegar a funcionar dentro de un área urbana cubriendo unos cuantos kilómetros para servir a miles de dispositivos finales; sin embargo, si los dispositivos finales necesitan enviar paquetes de información constantemente, la cantidad de dispositivos que un gateway puede manejar se reduce a unos cientos y su área de cobertura se vuelve desigual limitando la escalabilidad de una red LoRa según los experimentos realizados en \cite{Petajajarvi2017PerformanceCoverage} e igualmente otras limitaciones ilustradas en \cite{Adelantado2017UnderstandingLoRaWAN}. 

Además de las ventajas y desventajas que presentan las tecnologías IoT al momento de ser utilizadas para desarrollar cualquier tipo de aplicación, en estos dispositivos la seguridad es un aspecto que cobra especial importancia al estar centrados en redes inalámbricas, debido a que carecen de barreras físicas de acceso y la información viaja a través del aire, lo que hace más vulnerables estas redes a personas no deseadas que intenten utilizar dicha información a su gusto\cite{Chiu2006Seguridad802.11}. El no tener una política de seguridad de la información clara y definida, lleva inevitablemente al acceso no autorizado a una red informática o a los equipos que en ella se encuentran y puede ocasionar en la gran mayoría de los casos graves problemas. Por eso es que es necesario recurrir a la seguridad informática, ya que esta es la disciplina que se ocupa de diseñar las normas, procedimientos, métodos y técnicas, orientados a proveer condiciones seguras y confiables, para el procesamiento de datos en los sistemas informáticos\cite{AgudeloGonzalez2014ElBASC}.

En esta era tecnológica, las radiocomunicaciones se han vuelto imprescindibles en nuestras vidas diarias, desde los dispositivos personales como los teléfonos móviles, los auriculares inalámbricos, la difusión de señales de televisión, hasta los equipos de interconexión de redes de uso doméstico o de oficina. Por tal motivo, para que las personas indeseadas en una red inalámbrica no puedan tener acceso a la misma, los fabricantes de tecnologías inalámbricas se asociaron creando la Wi-Fi Alliance y formaron un estándar basado en las tecnologías cooperantes, así mismo la IEEE (Institute of Electrical and Electronics Engineers) hizo público el estándar 802.11 para las redes inalámbricas de área local\cite{RojasBautista2007SeguridadInalambricas}, de igual forma las redes ad hoc siguen técnicas patentadas o se basan en el estándar Bluetooth, que fue desarrollado por un consorcio de compañías comerciales que integran el Bluetooth Special Interest Group (SIG)\cite{Karygiannis2002WirelessSecurity}. 

Aunque la tecnología Wi-Fi es una de las tecnologías líder en las comunicaciones inalámbricas, hay un aspecto que en muchas ocasiones pasa inadvertido: la seguridad. Una de las vulnerabilidades que tiene el método de transmisión Wi-Fi es el tipo de protocolo de autenticación que se define en su utilización\cite{CorlettiEstrada2011SeguridadNiveles}, uno de los primeros tipos de autenticación que se utilizó en el protocolo, es la autenticación WEP (Wired Equivalent Privacy), dicho protocolo de encriptación está definido por defecto para redes WLAN y fue empleado para brindar protección a las redes inalámbricas al haber sido incluido en la primera versión del estándar IEEE 802.11 y mantenido sin cambios en versiones posteriores como los estándares 802.11a y 802.11b, con el fin  de garantizar compatibilidad de conexión entre distintos fabricantes\cite{Lehembre2006SeguridadWPA2}. Este sistema emplea el algoritmo RC4 para el cifrado de las llaves que pueden ser de 64 o 128 bit teóricos, puesto que en realidad son 40 o 104 y el resto (24bits) se emplea para el vector de inicialización. El protocolo WEP no fue creado por expertos en seguridad o criptografía, así que pronto se demostró que era vulnerable debido a los problemas presentes en el algoritmo RC4 que fueron descritos por David Wagner, por lo que en los años posteriores se definieron protocolos como WPA (Wireless Protected Access) y WPA2\cite{Lehembre2006SeguridadWPA2}, en donde, aunque se han descubierto algunas pequeñas debilidades desde sus lanzamientos, ninguna de ellas es peligrosa si se siguen unas mínimas recomendaciones de seguridad.  

Las tecnologías inalámbricas están cambiando rápidamente, nuevos productos y características se están introduciendo continuamente. Muchos de estos productos ahora ofrecen características de seguridad diseñadas para resolver debilidades de comunicación a larga distancia, envío de largas cadenas de información o abordar problemas recientemente descubiertos. Sin embargo, con cada nueva capacidad, es probable que surja una nueva amenaza o vulnerabilidad. Por lo tanto, es esencial estar al tanto de las tendencias actuales y emergentes en las tecnologías y en la seguridad o inseguridad de estas tecnologías. 

En el caso de LoRa los métodos de activación disponibles para los dispositivos proporcionan un alto nivel de seguridad mediante el cifrado simétrico en los intercambios de mensajes entre los nodos finales y los servidores. Los mecanismos proporcionados eliminan en gran medida la posibilidad de que un atacante inyecte paquetes maliciosos o suplante componentes de la red en el proceso de activación sin conocer las claves del nodo final. El servidor de la aplicación también es un componente sensible, ya que genera y almacena las claves utilizadas por los nodos finales en el proceso de activación y el cifrado de mensajes. Las claves deben almacenarse de forma segura y solo ser accesibles por el servidor de aplicación. El servidor de aplicación también emplea una interfaz de programación de aplicaciones (API). Por lo general, la API expone los datos por medio de un servidor web. Por lo tanto, el servidor web debe implementar una conexión segura utilizando Hyper Text Transfer Protocol Secure (HTTPS), basado en un certificado TLS firmado por una autoridad de certificación global (por ejemplo, Comodo o Let's Encrypt)\cite{Oniga2017AArchitecture}. Sin embargo, LoRa tiene inconvenientes inherentes causados por los compromisos que se han hecho en su diseño.

%Una conexión de red con dispositivos LoRa puede dividirse en dos partes fundamentales, una primera sección que va desde los nodos finales hasta el gateway y otra sección comprendida desde el gateway hasta los servidores. Dentro de la segunda sección existen diversas soluciones de seguridad a disposición ya que esta no es exclusiva para dispositivos LoRa, en cambio la primera sección puede ser susceptible a ataques de seguridad, por lo que esta porción de la red no puede considerarse una entidad de red confiable. 

Por ejemplo, existe una propuesta en donde la implementación de la %seguridad LoRa para los ataques se podría volver más firme usando %algoritmos de cifrados más robustos, utilizando protocolos de %autenticación y reglas de acceso al medio\cite{Aras2017}. De igual %forma, es viable encontrar más propuestas planteadas por otros %autores o dado el caso, ya que es una tecnología relativamente %nueva, generar otras soluciones al problema de seguridad presentado %en los dispositivos que utilizan el protocolo LoRaWAN. 

Además se pueden plantear otros temas de interés como el determinar %cómo implementar un esquema de seguridad de datos entre los nodos y %el gateway en una red LoRaWAN.
